% ==========================================================================
% preamble.tex — konfigurasi kelas, tata letak B5, dan paket
% Mesin kompilasi: XeTeX (via Tectonic). Kompatibel juga dengan pdfLaTeX.
% ==========================================================================

% --- Bahasa & enkoding -----------------------------------------------------
\usepackage{iftex}
\ifPDFTeX
  \usepackage[utf8]{inputenc}
  \usepackage[T1]{fontenc}
\fi
\usepackage[bahasa]{babel}   % istilah "Bab", "Daftar Isi", "Gambar", dsb.

% --- Tata letak halaman B5 (176 x 250 mm) ----------------------------------
\usepackage[
  b5paper,
  inner=22mm, outer=18mm,
  top=22mm,  bottom=24mm,
  headsep=8mm, footskip=12mm
]{geometry}

% --- Tipografi -------------------------------------------------------------
\usepackage{microtype}
\usepackage{csquotes}
\linespread{1.06}
\setlength{\parindent}{1.4em}
\usepackage{emptypage}       % halaman kosong benar-benar kosong

% Batasi kedalaman subbab: hanya Bab + Subbab (1 level) yang bernomor & masuk
% Daftar Isi. Judul "###" lama tetap tampil sebagai kepala tak bernomor.
\setcounter{secnumdepth}{1}
\setcounter{tocdepth}{1}

% --- Header / footer -------------------------------------------------------
\usepackage{fancyhdr}
\pagestyle{fancy}
\fancyhf{}
\fancyhead[LE]{\small\thepage\quad\leftmark}
\fancyhead[RO]{\small\rightmark\quad\thepage}
\renewcommand{\headrulewidth}{0.4pt}
\renewcommand{\chaptermark}[1]{\markboth{\thechapter.\ #1}{}}
\renewcommand{\sectionmark}[1]{\markright{\thesection.\ #1}}

% --- Matematika ------------------------------------------------------------
\usepackage{amsmath}
\usepackage{amssymb}
\usepackage{amsthm}

% --- Gambar & tabel --------------------------------------------------------
\usepackage{graphicx}
\graphicspath{{figures/}}
\usepackage{booktabs}
\usepackage{float}
\usepackage[font=small,labelfont=bf,skip=6pt]{caption}

% --- Warna -----------------------------------------------------------------
\usepackage{xcolor}
\definecolor{codebg}{RGB}{248,248,248}
\definecolor{codeframe}{RGB}{221,221,221}
\definecolor{kw}{RGB}{0,70,160}
\definecolor{str}{RGB}{176,0,64}
\definecolor{cmt}{RGB}{100,120,100}
\definecolor{lnum}{RGB}{150,150,150}

% --- Kode (listings) -------------------------------------------------------
\usepackage{listings}
\lstdefinestyle{pybook}{
  language=Python,
  backgroundcolor=\color{codebg},
  basicstyle=\ttfamily\footnotesize,
  keywordstyle=\color{kw}\bfseries,
  stringstyle=\color{str},
  commentstyle=\color{cmt}\itshape,
  numberstyle=\tiny\color{lnum},
  numbers=left, numbersep=8pt,
  frame=single, rulecolor=\color{codeframe},
  breaklines=true, breakatwhitespace=true,
  showstringspaces=false,
  tabsize=4,
  columns=fullflexible,
  keepspaces=true,
  upquote=true,
  aboveskip=10pt, belowskip=10pt,
  literate=%
    {á}{{\'a}}1 {é}{{\'e}}1 {í}{{\'i}}1 {ó}{{\'o}}1 {ú}{{\'u}}1
    {->}{{$\rightarrow$}}2
}
\lstset{style=pybook}

% --- Makro bantu bawaan pandoc (dibutuhkan fragmen .tex) -------------------
\providecommand{\tightlist}{\setlength{\itemsep}{0pt}\setlength{\parskip}{0pt}}
\providecommand{\passthrough}[1]{#1}
\providecommand{\pandocbounded}[1]{#1}
\providecommand{\real}[1]{#1}

% --- Tautan (hyperref selalu terakhir) -------------------------------------
\usepackage[hidelinks]{hyperref}
\hypersetup{
  bookmarksnumbered=true,
  pdftitle={Rekayasa Fitur Modern},
  pdfauthor={Fatma Indriani}
}
